\documentclass[10pt,a4paper]{amsart}
\usepackage[margin=19mm]{geometry}
\usepackage{amsmath,amssymb,mathtools}
\usepackage[scr=rsfs,cal=euler]{mathalfa}
\usepackage{xfrac}
\usepackage[unicode,hidelinks,bookmarks=false]{hyperref}
\newcommand{\ie}{i.e.\ }
\newcommand{\cf}{cf.\ }
\newcommand{\resp}{respectively\ }
\newcommand{\Subsection}{\subsection}
\providecommand{\nobreakdash}{\nobreak\hbox{-}}
\setlength{\emergencystretch}{2em}
\multlinegap0pt
\usepackage{bm}
\usepackage{bbm}
\usepackage{dsfont}
\usepackage{xspace}
\usepackage{cancel}
\usepackage{mathtools}
\usepackage{amsthm}
\makeatletter
\@ifundefined{lemma}{\newtheorem{lemma}{Lemma}}{}
\renewenvironment{proof}[1]{\par\medskip\noindent\textbf{#1}\enspace\ignorespaces}{\par\medskip}
\makeatother
\newcommand{\R}{ \mathbb{R} }
\newcommand{\tB}{\tilde{B}}
\newcommand{\tilb}{\tilde{b}}
\newcommand{\tlam}{\tilde{\lambda}}
\renewcommand{\tilde}{\widetilde}
\title[A Re-solving Heuristic for Dynamic Assortment Optimization w]{A Re-solving Heuristic for Dynamic Assortment Optimization with Knapsack Constraints}
\author{Xi Chen, Mo Liu, Yining Wang, Yuan Zhou}
\date{Source: arXiv:2407.05564v2 (CC BY 4.0)}
\begin{document}
\maketitle

\noindent\emph{Source and licence.} A Re-solving Heuristic for Dynamic Assortment Optimization with Knapsack Constraints. Version arXiv:2407.05564v2 (CC BY 4.0), distributed under the Creative Commons Attribution 4.0 International licence, as recorded in the arXiv licence metadata for this exact version and shown on the abstract page https://arxiv.org/abs/2407.05564v2. Full rights notice: licenses/arxiv-2407.05564-CC-BY-4.0.txt.\par\smallskip

\noindent\textit{Editorial scope.} Counted unit: the Lemma whose source label is \texttt{lem:hessian} in the article (source0.tex lines 1461--1465, source label \texttt{lem:hessian}), delivered together with the article's own complete original proof of it (source0.tex lines 1467--1516, one \texttt{proof} environment of 50 lines). No mathematics is written, repaired or re-derived: statement and proof are the article's own text line for line. Required context retained uncounted: none. Every definition, display and label of those lines is retained unchanged. Omitted from this package and not redistributed: 1--1460, 1466--1466, 1517--1596. Nothing inside a retained excerpt is omitted or altered: every retained body is delivered exactly as the article's lines are written, so no omission receipt of any kind accompanies this package. Citations: the retained excerpts carry no citation command at all, so no bibliography entry of the article is transcribed and no reference is redistributed. Reference adaptation: none was needed and none was made; every reference inside the delivered unit resolves inside it, so no unresolved reference or citation is produced. Printed slips kept literal: no printed word, symbol, hypothesis or number of the counted statement or of its proof was altered. Layout and class: the article's own class is \texttt{informs3aa} with packages including subfigure, arydshln, makecell, endnotes, longtable, tabularx, epstopdf, booktabs, etex, multirow, array, bbm, appendix, amsbsy; the wrapper is the lane's pinned \texttt{amsart} editorial wrapper at 10pt on A4, which restates the theorem-like environments the retained text needs and adds the article's own macro definitions that the retained text needs (\texttt{\textbackslash R}, \texttt{\textbackslash tB}, \texttt{\textbackslash tilb}, \texttt{\textbackslash tilde}, \texttt{\textbackslash tlam}), with extra wrapper packages (bm, bbm, dsfont, xspace, cancel, mathtools). The delivered numbering is the unit's own. Assets: no retained span contains a figure environment, a graphics-inclusion command, a PDF-page inclusion, a table or a TikZ picture; no image file is copied into this package, the package declares no asset dependency and no image file is redistributed with it. Licence: version arXiv:2407.05564v2 of this article is distributed under the Creative Commons Attribution 4.0 International licence, as recorded in the arXiv licence metadata for this exact version and shown on the abstract page https://arxiv.org/abs/2407.05564v2. A full rights notice is delivered at licenses/arxiv-2407.05564-CC-BY-4.0.txt. Characters: every retained excerpt is pure ASCII, so the delivered engine, XeLaTeX with its default Latin Modern text font, reports no missing character.
\begin{lemma}
	Under Assumptions (A1) to (A3), for any $\gamma,s$ satisfying $\|\gamma-\gamma_0\|_{\infty}\leq\rho_0$ and $|s-s^{\tau_0}|\leq\rho_0$, $\Psi$ is differentiable in $(\gamma,s)$ and furthermore $
	\|\partial^2 \Psi(\gamma,s)\|_{\mathrm{op}}\leq \frac{8 M^2 KN B_0^3}{\chi_0^2}.$
	\label{lem:hessian}
\end{lemma}

\begin{proof}{Proof.}
	Recall that $\Psi$ can be written as the following linear programming problem:
	\begin{align*}
		\Psi(\gamma,s) = & \max_{x\in [0,1]^N} \sum_{i=1}^N r_i v_i  x_i,s.t.\sum_{i=1}^Nx_i \leq K,  1+\sum_{i=1}^Nv_ix_i\leq s, \sum_{i=1}^N A_{ji}v_ix_i\leq \gamma_j\left[1+\sum_{k=1}^Nv_kx_k\right],\;\forall j\in[M].
	\end{align*}
	We convert the constraints of $\Psi$ into the standard form $\tilde{A} x \le \tilde b$, where $\tilde{A} \in \R^{(M+2+ 2N)\times N}$ denotes the coefficient matrix in the left hand side of constraints, and $\tilb \in \R^{M+2+ 2N}$ denotes the vector in the right hand side of constraints. The first row of matrix $\tilde{A}$ is an all-one vector. For row index $i$  from 2 to $M + 1$, each entry in $\tilde{A}$ is $(A_{(i-1), j} - \gamma_{i - 1} ) v_{j} $, for column index $j = 1,...,N$. The $(M+2)$-th row of $\tilde{A}$ is the same as $v$. The last $2N\times N$ submatrix of $\tilde{A}$ is $[\mathbb{I}_N, -\mathbb{I}_N]^\top$.  The vector $\tilb$ is $[K, \gamma, s-1, [1,...,1]_N, [0,...,0]_N]^\top$.
	We use $\tlam^* \in \R^{M+2+ 2N}$ to denote the optimal dual variables of constraints $\tilde{A} x \le \tilde b$ for $\Psi$.  Notice that both $\tlam^*$ and $x^*$ depend on $\gamma$ and $s$. 
	
	We first consider the first order derivative of $\Psi(\gamma, s)$ using the envelope theorem. We notice that $\gamma_j$ only appears in the $(j+1)$-th row in constraints, and $s$ only appears in the last row of constraints. Thus, according to the envelope theorem, the partial derivatives of \(\Psi\) with respect to the parameters \(\gamma\) and \(s\) are given by:
	$ \frac{\partial \Psi}{\partial \gamma_j} = \tlam_{j+1}^* \left( \frac{\partial \tilde b}{\partial \gamma_j} - \frac{\partial \tilde{A} }{\partial \gamma_j} x^* \right)_{(j+1)}, $ and 
	$ \frac{\partial \Psi}{\partial s} = \tlam_{M+2}^* \left( \frac{\partial\tilde  b}{\partial s} - \frac{\partial \tilde{A}}{\partial s} x^* \right)_{(M+2)}. $
	
	By the value of $\tilde{A}$ and $\tilde b$, we have that $  \left( \frac{\partial \tilde b}{\partial \gamma_j} - \frac{\partial \tilde{A} }{\partial \gamma_j} x^* \right)_{(j+1)} = 1 + \sum_{i} r_i v_i x_i$, and $ \left( \frac{\partial  \tilde b}{\partial s} - \frac{\partial \tilde{A}}{\partial s} x^* \right)_{(M+2)} = 1$. Thus, we obtain that $
		\frac{\partial \Psi}{\partial \gamma_j} =  \tlam_{j+1}^*(1 +\sum_{i} r_i v_i x_i), \text{for } j \in [1,....,M], \text{and } ~~\frac{\partial \Psi}{\partial s} = \tlam_{M+2}^* $. 
	
	
	
	
	
	
	Next, we consider the second derivative of $\Psi(\gamma, s)$. We denote the basis of the binding constraints at $x^*$ by $\tB^*(\gamma, s) \in \R^{N\times N}$ and denote its corresponding right hand side by $\tilb_B \in \R^{N}$. Thus, we have that $\tB^*(\gamma, s) x^* = \tilb_B$. We further denote the optimal dual variables for the binding constraints by $\tlam_B \in \R^N$. By the complementary slackness, the optimal dual variables for nonbinding constraints are zero, so we have that $\tB^*(\gamma, s)^\top \tlam_B = v \odot r$, where $\odot$ represents the entrywise product of two vectors.
	
	The major part of anlayzing the hessian matrix $\partial^2 \Psi(\gamma,s)$ is to examine how $\tlam^*$ and $x^*$ depend on $(\gamma, s)$. By Assumption (A3), for any $(\gamma,s) \in \{(\gamma,s): \|\gamma - \gamma_0\| \le \rho_0 \wedge |s-s^{\tau_0}|\leq\rho_0  \}$,  the minimum singluar value of $\tB^*$ is greater than or equal to $\chi_0>0$. It implies that there exsits a unique solution $x^*$ for $\Psi(\gamma, s)$, and that the basis $\tB^*$ remains unchanged when $(\gamma, s)$ is within the set of $\{(\gamma,s): \|\gamma - \gamma_0\| \le \rho_0 \wedge |s-s^{\tau_0}|\leq\rho_0  \}$. Thus, for the non-binding constraints regarding to resource $j$, for any $k = 1,...,M$, we have that $ \frac{\partial \tlam^*_{j+1}}{ \partial \gamma_k }  = 0, ~~ \frac{\partial \tlam^*_{j+1}}{ \partial s }  = 0.$
	Similarly, for the non-binding constraints regarding to resource $j$, we have that 
	$ \frac{\partial x^*}{ \partial \gamma_j }  = 0.$
	
	Thus, when analyzing the upper bound for the operator norm of the hessian matrix, $\|\partial^2 \Psi(\gamma,s)\|_{\mathrm{op}}$, it suffices to consider the binding constraints. Since $\tB^*(\gamma, s)^\top \tlam_B = v$ and the minimum singular value of $\tB^*$ is no smaller than $\chi_0$,  we have that $\|v\| = \| \tB^*\tlam_B \| \ge\chi_0  \|\tlam_B\|$. Since $\|v\| \le \sqrt{N}B_0$, we have that $\|\tlam_B\| \le \frac{ \sqrt{N}B_0}{\chi_0}$.
	By taking the derivative of $(\gamma, s)$ on both sides of $\tB^*(\gamma, s)^\top \tlam_B = v\odot r$ , we obtain that $\frac{\partial \tB^*}{ \partial (\gamma, s) } \tlam_B +  \tB^*  \frac{\partial \tlam_B}{ \partial (\gamma, s) } = 0.$ Therefore, we have that $
		\left\| \tB^*  \frac{\partial \tlam_B}{ \partial (\gamma, s) } \right\| = \left\| \frac{\partial \tB^*}{ \partial (\gamma, s) } \tlam_B \right\|.$ 
	
	Since  $\tB$ is a submatrix of $\tilde{A}$, we have that each entry of $ \frac{\partial \tB^*}{ \partial (\gamma, s) }$ is smaller than or equal to $B_0$. Thus, the right hand side $\left\| \frac{\partial \tB^*}{ \partial (\gamma, s) } \tlam_B \right\| $ is smaller than or equal to $ (M+1 ) B_0 \times   \frac{ \sqrt{N}B_0}{\chi_0} =  \frac{ (M+1)\sqrt{N}B_0^2}{\chi_0}$. Thus, we have that $\left\| \tB^*  \frac{\partial \tlam_B}{ \partial (\gamma, s) } \right\| \le \frac{ (M+1)\sqrt{N}B_0^2}{\chi_0}$. 
	Again, by the minimum singular value of $\tB^*$, we have that 
	\begin{align}\label{equ2:lemma}
		\left \|   \frac{\partial \tlam_B}{ \partial (\gamma, s) } \right\| \le \frac{ (M+1)\sqrt{N}B_0^2}{\chi_0^2}.
	\end{align}
	
	Next, we consider the upper bound for $\left \|   \frac{\partial  x^*}{ \partial (\gamma, s) } \right\|$. Since $\tB^* x^* = \tilb_B$, by taking the derivative on both sides, we have that  $\frac{\partial \tB^*}{ \partial (\gamma, s) } x^* +  \tB^*  \frac{\partial x^*}{ \partial (\gamma, s) } = \frac{\partial \tilb_B}{ \partial (\gamma, s) }. $
	
	
	Since $\tilb_B$ is subvector of $\tilb$, we have that $\| \frac{\partial \tilb_B}{ \partial (\gamma, s) }\| \le \sqrt{N}$.
	By $\|x^*\|_2 \le \|x^*\|_1 \le K$, we have that  $\left\| \tB^*  \frac{\partial x^*}{ \partial (\gamma, s) } \right\| \le \sqrt{N} + (M+1)B_0 K$. Again, by the minimum singular value of $\tB^*$, we have that
	
	\begin{align}\label{equ3:lemma}
		\left \|   \frac{\partial x^*}{ \partial (\gamma, s) } \right\| \le \frac{\sqrt{N} + (M+1)B_0 K}{\chi_0}.
	\end{align}
	
	By taking the derivatives of $\frac{\partial \Psi}{\partial \gamma}$ and $\frac{\partial \Psi}{\partial s}$, we have that for the hessian matrix $\partial^2 \Psi(\gamma,s)$, each entry is no larger than $ \left \|   \frac{\partial \tlam_B}{ \partial (\gamma, s) } \right\|  (1 + B_0 K) + \|\tlam_B\| \cdot B_0 \cdot \left \|   \frac{\partial x^*}{ \partial (\gamma, s) } \right\|. $
	Combining the upper bounds in \eqref{equ2:lemma} and \eqref{equ3:lemma}, we have that each entry in the hessian matrix is no larger than $ \frac{ (M+1)\sqrt{N}B_0^2 (1 + B_0 K)}{\chi_0^2} +  \frac{ \sqrt{N}B_0^2}{\chi_0}   \frac{\sqrt{N} + (M+1)B_0 K}{\chi_0},$
	which is further no larger than $\frac{4 MKN B_0^3}{\chi_0^2}$. Thus, the operator norm of the hessian matrix $\partial^2 \Psi(\gamma,s)$ is no larger than $\frac{8 M^2 KN B_0^3}{\chi_0^2}$.
\end{proof}

\end{document}
